% ============================================================
% Validation — tests on mocks, simulations, or objects with
% known ground truth. Astronomy has no held-out benchmark
% labels, so this section is where you earn the reader's trust
% in the method's accuracy and error bars. (Rename freely:
% "Robustness tests", "Simulations", "Mock tests".)
% ============================================================
\section{Validation}
\label{sec:validation}

\subsection{Mock Construction}
\label{subsec:mocks}

% Build simulated data with known ground truth that matches
% the REAL data's noise, resolution, and selection function —
% state explicitly how realism was achieved (e.g. injecting
% models into real images rather than pristine simulations).

\subsection{Recovery Tests}
\label{subsec:recovery}

% How well does the method recover the input parameters?
% Quantify biases and scatter vs. signal-to-noise, size,
% complexity. A scatter plot of recovered vs. input values is
% the canonical figure:
% \begin{figure}[htbp]
%     \centering
%     \includegraphics[width=\linewidth]{figures/recovery.pdf}
%     \caption{Recovered vs. input parameters from mock tests.
%              The dashed line shows the one-to-one relation.}
%     \label{fig:recovery}
% \end{figure}

\subsection{Comparison with Established Methods}
\label{subsec:comparison}

% Where alternatives exist (cite them, e.g.
% \citealp{2002AJ....124..266P}), compare on shared objects or
% mocks and explain — not just tabulate — the differences.
